\documentclass[11pt,letterpaper,english]{article}
\usepackage{microtype}
\usepackage{titling}
\usepackage{amsmath} 
\usepackage{amsthm}
\usepackage{mathrsfs}
\usepackage{caption}
\usepackage{bibunits}
\usepackage{fontawesome}
\usepackage[font=small,labelfont=bf]{caption}
\usepackage{subcaption}
\usepackage{setspace}
\usepackage{amssymb}
\usepackage{multirow}
\usepackage{rotating}
\usepackage{array}
\usepackage{multibib}
\usepackage{blindtext}
\usepackage{mathtools}
\usepackage{siunitx}
\usepackage{colortbl}
\usepackage{float}
\usepackage[flushleft]{threeparttable}
\usepackage{algorithm}
\usepackage{multirow, makecell}
\makeatletter
\renewcommand{\ALG@name}{Pseudo-code}
\makeatother
\usepackage{algorithmic}
\usepackage{xr}
\usepackage{color}
%\externaldocument{appendix_jan1}
\DeclareMathAlphabet{\mathcalligra}{T1}{calligra}{m}{n}
\newcommand{\Var}{\mathrm{Var}}
\DeclareMathOperator{\E}{\mathbb{E}}
\DeclareMathOperator{\Tr}{Tr}
\usepackage[sectionbib,round]{natbib}
\setlength{\bibsep}{0.0pt}
\setlength{\footskip}{30pt}
\newtheorem{propn}{Proposition}\newtheorem{lemm}{Lemma}
\DeclareMathOperator*{\argmax}{arg\,max}
\DeclareMathOperator*{\argmin}{arg\,min}
\onehalfspacing
\usepackage{times}
\usepackage[small,compact]{titlesec} 
\usepackage{abstract}
\renewcommand{\abstractnamefont}{\normalfont\bfseries}
\renewcommand{\abstracttextfont}{\normalfont}
\usepackage[hang,flushmargin]{footmisc}
\usepackage{fullpage} 
\usepackage{hyperref}
\usepackage[compact]{titlesec}
\usepackage[section]{placeins}
\usepackage{footnote}
\usepackage{epstopdf}
\usepackage{placeins}
\usepackage[toc,page]{appendix}
\usepackage{blkarray}
\usepackage{stackengine}
\usepackage{accents}
\usepackage{enumitem}
\usepackage[mathscr]{euscript}
\usepackage{booktabs}  
\usepackage{nicefrac}
\usepackage{longtable}

\newcommand\norm[1]{\left\lVert#1\right\rVert}

\newtheorem{definition}{Definition}[section]
\newcommand*\mean[1]{\bar{#1}}
%\theoremstyle{definition}
%\newtheorem{defn}{Definition}

\usepackage[utf8]{inputenc}
\usepackage{longtable}
\usepackage{graphicx}
\usepackage{epstopdf}
\DeclareGraphicsRule{.tif}{png}{.png}{`convert #1 `dirname #1`/`basename #1 .tif`.png}
\usepackage{bbm}
\usepackage{dsfont}
\usepackage{comment}
\usepackage{pgfplots}

\newcommand{\squishlist}{
   \begin{list}{$\bullet$}
    { \setlength{\itemsep}{0pt} \setlength{\parsep}{1pt}
      \setlength{\topsep}{1pt} \setlength{\partopsep}{1pt}
      \setlength{\leftmargin}{1.5em} \setlength{\labelwidth}{1em}
      \setlength{\labelsep}{0.5em} } }

\newcommand{\squishlisttwo}{
   \begin{list}{$\bullet$}
    { \setlength{\itemsep}{0pt} \setlength{\parsep}{0pt}
      \setlength{\topsep}{0pt} \setlength{\partopsep}{0pt}
      \setlength{\leftmargin}{1em} \setlength{\labelwidth}{1.5em}
      \setlength{\labelsep}{0.5em} } }

\theoremstyle{definition}
\newtheorem{assum}{Assumption}

\newcommand{\todo}[1]{\textcolor{blue}{#1}}

\theoremstyle{definition}
\newtheorem{defn}{Definition}

\newcommand{\squishend}{
    \end{list}  }

\newcommand{\ebi}[1]{\noindent{\textcolor{blue}{\{{\bf Ebrahim:}  #1\}}}}
\newcommand{\hema}[1]{\noindent{\textcolor{green}{\{{\bf Hema:}  #1\}}}}

\defaultbibliographystyle{abbrvnat} 
\defaultbibliography{ref} 

\begin{document}
\title{Embedded Auditing in Experimentation-Driven Digital Platforms}


\author{
       Ebrahim Barzegary \\ %\thanks{}  \\
        ESSEC Business school\\
        %\and
       %Hema Yoganarasimhan\footnotemark[1]\\
        %University of Washington
 }

\thispagestyle{empty}

\date{}
\maketitle

\begin{abstract}

Digital platforms increasingly rely on large-scale, continuous experimentation to optimize user-facing interventions such as pricing, recommendations, and persuasive content. While experimentation-driven systems improve learning and efficiency, they create a governance challenge: optimized interventions may generate heterogeneous and localized harms that remain obscured by positive average outcomes. Risk-based regulatory frameworks such as the European Union’s Digital Services Act require ongoing assessment and mitigation of such risks, yet operational auditing mechanisms for continuously optimized systems remain underdeveloped.

We introduce \emph{embedded auditing}, a governance framework that leverages the randomized experimentation infrastructure already used for optimization. Extending recursive partitioning to a two-outcome setting, we identify interpretable user segments based on joint treatment effects on a firm-side and a consumer-side outcome. The method surfaces segments where firm gains coincide with consumer losses, transforming heterogeneous outcome divergence into transparent, rule-based diagnostics without requiring access to proprietary models.

By integrating auditing directly into experimentation pipelines, the framework operationalizes risk-based governance for digital platforms and demonstrates how optimization infrastructure can serve as a mechanism of accountability.

\vspace{5mm}
\noindent \textbf{Keywords:} algorithmic governance, embedded auditing, experimentation-driven systems, heterogeneous effects, digital platforms, accountability
\end{abstract}


\begin{bibunit}
\newpage

\section{Introduction}
\label{sec:intro}

Digital platforms increasingly operate as experimentation-driven systems. Prices, recommendation policies, interface designs, promotional messages, and subscription terms are routinely refined through large-scale, continuous A/B testing optimized against firm-side performance metrics such as conversion, engagement, and revenue \citep{kohavi2020trustworthy}. Continuous experimentation has evolved into an organizational capability supported by dedicated infrastructure, metric governance systems, and institutionalized decision routines \citep{auer2021_continuous_experimentation, quinn2024_ab_slr, dmitriev2021_ab_flywheel}. Advances in machine learning and generative AI further accelerate this process by enabling rapid generation and iterative testing of personalized content at scale \citep{ye2025lola}. In many digital markets, experimentation is not an auxiliary tool but a central mechanism of strategic learning and competitive advantage.

At the same time, experimentation-driven optimization creates a governance challenge. Interventions are typically evaluated using aggregate firm-level metrics, while consumer-side consequences—such as satisfaction, regret, stress, overuse, or long-term retention—may receive secondary weight unless they directly affect financial performance. As a result, interventions can produce positive average outcomes while imposing localized costs on specific consumer segments. Research in marketing, law, and information systems has begun to document such algorithmic harms, including cases where optimization processes systematically disadvantage particular groups or amplify behavioral vulnerabilities \citep{bar2023algorithmic, lambrecht2019algorithmic}. Crucially, these harms are often heterogeneous rather than uniform: the same intervention may benefit some users while disadvantaging others.

Heterogeneity poses a fundamental obstacle for auditing and governance. Risk-based regulatory frameworks such as the European Union’s Digital Services Act (DSA) require designated platforms to identify, assess, and mitigate systemic risks arising from the design and functioning of their services, including algorithmic systems \citep{dsa2022article34}. Similarly, lifecycle-oriented governance regimes emphasize ongoing monitoring rather than one-time compliance. Yet regulatory language typically presumes that harms can be defined and detected at a system level. In experimentation-driven environments, harms may be segment-specific and statistically obscured by positive average effects. When outcome divergence is localized, aggregate evaluation can mask precisely the cases that governance seeks to identify. While regulatory frameworks articulate recurring risk assessment duties, they provide limited operational guidance on how heterogeneous consumer harms should be measured and monitored within high-frequency experimentation environments.

The broader literature on algorithmic accountability reinforces these concerns. Foundational scholarship argues that meaningful accountability requires institutional mechanisms that connect design choices, implementation processes, and observable outcomes \citep{kroll2017_accountable_algorithms}. Internal algorithmic auditing frameworks emphasize lifecycle documentation and post-deployment monitoring rather than static compliance checks \citep{raji2020_internal_audit}. At the same time, auditing regimes face structural tensions: external audits often depend on access to proprietary systems, raising confidentiality and trade-secret concerns, and audit processes can be vulnerable to “audit-washing,” whereby firms shape audit scope or selectively disclose favorable evidence without addressing localized harms \citep{goodman2022_auditwashing}. In fast-moving digital markets characterized by continuous experimentation, oversight mechanisms that operate episodically or rely on aggregate reporting are likely to lag behind system evolution.

This paper proposes an alternative governance mechanism: \emph{embedded auditing}. Rather than auditing algorithms from outside the optimization process, we leverage the randomized experimentation infrastructure firms already use. Randomized A/B tests generate causal variation by design and therefore contain the information necessary to evaluate heterogeneous effects of interventions across user segments. By integrating auditing directly into experimentation pipelines, oversight becomes continuous rather than episodic and does not require disclosure of proprietary model architecture. In this sense, experimentation infrastructure—originally developed for optimization—can also function as an evidence-generating backbone for governance.

Methodologically, we extend recursive partitioning \citep{athey2016recursive} to a two-outcome setting in which treatment effects on a firm-side outcome (e.g., revenue) and a consumer-side outcome (e.g., usage, satisfaction, or regret) are jointly evaluated. The algorithm identifies interpretable segments defined by pre-treatment characteristics and prioritizes splits that reveal meaningful heterogeneity in both outcomes. A configurable region-weighting component allows auditors to emphasize segments where firm gains coincide with consumer losses—cases that are particularly salient from a governance perspective. By transforming heterogeneous outcome divergence into transparent, rule-based segment diagnostics, the framework provides a reproducible audit primitive tailored to experimentation-driven systems.

The proposed approach addresses three practical challenges in governing adaptive digital platforms. First, it operationalizes risk-based governance using experimentally identified treatment effects rather than abstract welfare concepts, narrowing the metric problem to observable, auditable quantities. Second, it makes heterogeneity central to oversight, ensuring that localized harms are not averaged away. Third, it reduces discretion and opportunities for audit-washing by relying on interpretable segmentation rules, explicit objective definitions, and a small number of cross-validated hyperparameters. Because auditing is embedded within existing experimentation workflows, it can operate at the same cadence as optimization without requiring access to proprietary models or trade-secret disclosure.

More broadly, the paper advances an institutional insight: experimentation infrastructures can serve not only as engines of optimization but also as mechanisms of accountability. In experimentation-driven digital platforms, effective governance requires continuous, segment-aware diagnostics aligned with the tempo of system updates. Embedded auditing provides one such mechanism, demonstrating how heterogeneous firm–consumer trade-offs can be rendered visible and actionable within existing organizational processes.


\section{Related Literature}
\label{sec:related_lit}

Digital platforms are increasingly organized around continuous experimentation. Large-scale A/B testing infrastructures now shape pricing, ranking, recommendation, interface design, and subscription policies \citep{kohavi2020trustworthy, auer2021_continuous_experimentation, quinn2024_ab_slr}. Rather than serving as episodic evaluation tools, experimentation systems have become embedded organizational capabilities that enable rapid deployment, measurement, and refinement of interventions. In digital marketing contexts, this infrastructure supports fine-grained personalization and scalable optimization of firm-side metrics such as conversion, engagement, and revenue.

Recent advances in machine learning and generative AI further intensify this dynamic by reducing the marginal cost of generating and iterating persuasive content \citep{ye2025lola}. As experimentation cycles shorten and adaptive learning systems operate continuously, optimization becomes tightly coupled with platform operation. What is measured and optimized becomes systematically amplified.

A substantial body of research in IS and marketing examines experimentation as 
an engine of organizational learning and performance improvement 
\citep{kohavi2020trustworthy, bojinnov2022_review}. These systems enable 
firms to tailor interventions to heterogeneous users and to identify superior 
policies efficiently. However, experimentation infrastructures are typically 
optimized with respect to narrow firm-side outcomes. Consumer-side consequences—
such as satisfaction, regret, retention quality, or longer-term well-being—are 
often secondary, delayed, or excluded from objective functions. When treatment 
effects are heterogeneous, aggregate improvements may mask localized losses.

Parallel research in IS addresses algorithmic accountability and governance 
\citep{kroll2017_accountable_algorithms, raji2020_internal_audit, rai2020explainable}. 
A recurring theme is that adaptive systems evolve continuously, rendering static 
audits insufficient. Governance mechanisms must operate at the same cadence as 
optimization. Risk-based regulatory frameworks similarly emphasize ongoing 
monitoring and systemic risk assessment in digital platforms 
\citep{dsa2022article34, eu2024aiact}. 

Yet practical auditing approaches face institutional and technical constraints. 
Model-based audits often require access to proprietary systems, while transparency 
mechanisms may not scale in complex adaptive environments. Outcome-based auditing 
offers an alternative by evaluating realized system behavior rather than internal 
architecture. However, most existing outcome-based approaches focus on aggregate 
disparities or static deployments and do not directly address segment-level 
divergence under continuous experimentation \citep{goodman2022_auditwashing}.

Heterogeneity complicates this problem. Marketing and economics research has 
documented that personalization can generate heterogeneous welfare effects in 
pricing, ranking, and advertising contexts \citep{dube2023personalized, 
ursu2018power, donnelly2024_welfare}. IS research has shown that algorithmic 
optimization may produce unintended disparities even without explicit demographic 
targeting \citep{lambrecht2019algorithmic}. Legal scholarship similarly highlights 
that algorithmic harms in consumer markets are often diffuse and concentrated 
in specific subpopulations \citep{bar2023algorithmic}. 

In experimentation-driven systems, such heterogeneity is amplified by adaptive 
learning. Policies selected for average improvement may systematically disadvantage 
particular segments. A positive average treatment effect does not preclude negative 
effects for identifiable users. Detecting such divergence requires moving beyond 
average effects toward structured heterogeneity diagnostics.

Methodologically, heterogeneous treatment effect estimation provides tools for 
identifying variation in causal effects across subpopulations 
\citep{athey2016recursive, wager2018estimation, athey2019grf}. Recursive 
partitioning approaches are especially attractive in organizational settings 
because they yield interpretable segmentation rules. However, most existing 
HTE methods are designed for a single outcome. Governance concerns in 
experimentation-driven systems arise precisely when treatment effects diverge 
across outcomes—for example, when firm-side gains coincide with consumer-side 
losses.

Taken together, prior research establishes that digital platforms rely on 
continuous experimentation for optimization, that adaptive algorithmic systems 
require ongoing governance, and that treatment effects are often heterogeneous. 
Yet experimentation infrastructures are designed to detect average performance 
improvements, while governance frameworks presume the ability to identify systemic 
risk. When harms are segment-specific and masked by aggregate gains, this mismatch 
becomes consequential. Integrating multi-outcome heterogeneous treatment effect 
analysis directly into experimentation pipelines provides a mechanism for aligning 
optimization infrastructure with governance requirements in adaptive digital 
marketing systems.

\section{Modeling Framework and Problem Setup}
\label{sec:app}

\subsection{Marketing Intervention and Post-Treatment Behavior}

We consider a common marketing setting in which firms conduct randomized experiments (e.g., A/B tests) to evaluate alternative marketing interventions. Let $T_i \in \{0,1\}$ denote the treatment assigned to individual $i$, where $T_i=0$ represents the control condition and $T_i=1$ represents the active intervention (e.g., a short free trial, an urgency message, or a personalized promotion). Each individual is characterized by a vector of observed pre-treatment covariates $X_i \in \mathcal{X}$, such as demographics, browsing behavior, or prior usage history.

For each individual, we observe two post-treatment outcomes:
\begin{itemize}
    \item A \emph{firm-side outcome} $Y_i^F$, capturing firm-relevant performance metrics such as conversion, subscription, purchase behavior, or revenue.
    \item A \emph{consumer-side outcome} $Y_i^C$, capturing post-treatment consumer outcomes such as product usage, satisfaction, retention, returns, or well-being.
\end{itemize}

The outcomes $Y_i^F$ and $Y_i^C$ may be binary, continuous, or discrete, and the framework does not impose distributional restrictions on their support. In many applications, both outcomes are directly observed for all individuals. In other settings, consumer-side outcomes may be partially observed due to post-treatment behavior (e.g., satisfaction observed only for retained users). The proposed method accommodates both cases and allows for estimation of consumer-side effects under conditional observability, provided sufficient information is available within segments.

We adopt the potential outcomes framework. Let $Y_i^F(t)$ and $Y_i^C(t)$ denote the potential firm-side and consumer-side outcomes for individual $i$ under treatment $t \in \{0,1\}$. Due to the fundamental problem of causal inference, we observe only the outcomes corresponding to the realized treatment $T_i$. We assume treatment is randomly assigned:
\[
T_i \perp\!\!\!\perp \{Y_i^F(0), Y_i^F(1), Y_i^C(0), Y_i^C(1)\} \mid X_i.
\]

\subsection{Objective and Formal Problem Definition}
\label{ssec:setting}

Our objective is not to optimize treatment assignment, but to \emph{audit} firm-optimized interventions by identifying subpopulations that exhibit systematically different joint responses in firm-side and consumer-side outcomes. Marketing interventions often generate heterogeneous effects across individuals, and the same policy may produce mutually beneficial outcomes for some subpopulations while generating trade-offs or harm for others. The goal of the framework is to uncover such heterogeneity and make these patterns explicit.

Let $\tau^F(x)$ and $\tau^C(x)$ denote the conditional average treatment effects for individuals with covariates $X_i = x$:
\[
\tau^F(x) = \mathbb{E}[Y_i^F(1) - Y_i^F(0) \mid X_i = x], \qquad
\tau^C(x) = \mathbb{E}[Y_i^C(1) - Y_i^C(0) \mid X_i = x].
\]

In some applications, consumer-side outcomes may only be observable for a subset of individuals (e.g., conditional on retention or engagement). In such cases, $\tau^C(x)$ is interpreted as a conditional treatment effect and is estimated using the corresponding observable subpopulation. The proposed method accommodates both fully observed and partially observed consumer outcomes by enforcing minimum support conditions within segments to ensure stable estimation.

Each subpopulation can be characterized by the sign and magnitude of $(\tau^F(x), \tau^C(x))$, giving rise to qualitatively different joint outcome patterns. For example, a segment with $\tau^F(x) > 0$ and $\tau^C(x) < 0$ corresponds to a setting in which the intervention improves firm-side performance while reducing consumer outcomes, whereas a segment with both effects positive reflects a mutually beneficial intervention. Other segments may experience uniformly negative effects or exhibit weaker responses on one dimension relative to the other.

Let $\mathcal{D} = \{(X_i, T_i, Y_i^F, Y_i^C)\}_{i=1}^n$ denote data from a randomized experiment. Our goal is to partition the covariate space $\mathcal{X}$ into $L$ disjoint and interpretable segments $\mathcal{S} = \{S_1, \dots, S_L\}$ such that treatment effects are approximately homogeneous within each segment. For a given segment $S_\ell$, we define:
\[
\tau^F(S_\ell) = \mathbb{E}[Y^F(1) - Y^F(0) \mid X \in S_\ell], \qquad
\tau^C(S_\ell) = \mathbb{E}[Y^C(1) - Y^C(0) \mid X \in S_\ell].
\]

The objective of the algorithm is to construct segmentations that reveal meaningful variation in joint firm--consumer treatment effects across segments while maintaining stability within segments. By making these joint outcome patterns explicit, the framework enables firms, researchers, and regulators to systematically audit how optimized interventions distribute benefits and costs across different subpopulations, rather than relying solely on aggregate treatment effects.


\section{Algorithmic Framework: Recursive Partitioning for Joint Outcome Analysis}

This section introduces the algorithmic framework that operationalizes the auditing objective defined in the modeling section. The goal of the algorithm is to construct interpretable customer segments in which marketing interventions exhibit heterogeneous \emph{joint} effects on firm-side and consumer-side outcomes, and to prioritize segments corresponding to outcome patterns of substantive interest.

\subsection{Segmentation via Recursive Partitioning}
\label{ssec:recur_part}

Recursive partitioning is a general-purpose meta-algorithm for dividing a covariate space into disjoint, interpretable segments by recursively applying axis-aligned splits. At each iteration, the algorithm considers candidate splits that divide an existing segment along a threshold of a single covariate, evaluates their contribution to a predefined objective function, and selects the split that yields the largest improvement. The process continues until a stopping criterion is met.

Formally, the procedure operates as follows:
\begin{itemize}
    \item \emph{Input}: An objective function $\mathcal{F}(\Pi)$ that evaluates a partitioning $\Pi$ of the covariate space.
    \item \emph{Initialization}: Start with an initial partition $\Pi_0$ consisting of the entire sample.
    \item \emph{Iteration}: For iteration $r=0,1,\dots$,
    \begin{itemize}
        \item For each current segment $\pi \in \Pi_r$, for each covariate $x \in \mathcal{X}$, and for each candidate threshold $v$ in the support of $x$ within $\pi$, compute the gain
        \[
        \Delta(\pi,x,v) = \mathcal{F}(\Pi_r + \{\pi,x,v\}) - \mathcal{F}(\Pi_r).
        \]
        \item Select the split $(\pi^*,x^*,v^*)$ that maximizes $\Delta(\pi,x,v)$.
        \item Update the partition $\Pi_{r+1} = \Pi_r + \{\pi^*,x^*,v^*\}$.
    \end{itemize}
\end{itemize}

The index $r$ denotes the iteration number, and $\Pi_r$ represents the partition after $r$ splits.

\subsection{Joint Splitting Criterion}
\label{sec:spliting_criterion}

The central contribution of the algorithm lies in the construction of a splitting criterion that jointly evaluates heterogeneity in firm-side and consumer-side treatment effects. Standard heterogeneous treatment effect methods typically optimize splits with respect to a single outcome. In contrast, our objective is to identify splits that simultaneously differentiate treatment effects on both outcomes, thereby revealing segments with distinct joint outcome patterns.

\vspace{5mm}
\noindent
\emph{Segment-level treatment effect estimation.}

For any segment $S \subset \mathcal{X}$, we estimate segment-level treatment effects using training data:
\[
\hat{\tau}^F(S) = \frac{1}{|S|} \sum_{i \in S} \big[ \hat{Y}^F_i(1) - \hat{Y}^F_i(0) \big], \qquad
\hat{\tau}^C(S) = \frac{1}{|S_C|} \sum_{i \in S_C} \big[ \hat{Y}^C_i(1) - \hat{Y}^C_i(0) \big],
\]
where $S_C$ denotes the subset of observations in $S$ for which the consumer-side outcome is observed. The estimators $\hat{Y}^F_i(t)$ and $\hat{Y}^C_i(t)$ are computed as within-segment sample means under treatment $t \in \{0,1\}$, exploiting random assignment for unbiased estimation.

Given a candidate split of a segment $S$ into two child segments $S_L$ and $S_R$, the gain from the split is defined as
\[
\Delta(S_L,S_R)
=
\sum_{j \in \{L,R\}} w_j \cdot g\!\left( \hat{\tau}^F(S_j), \hat{\tau}^C(S_j) \right),
\]
where $w_j = |S_j|/|S|$ and $g(\cdot,\cdot)$ is a joint heterogeneity function.

We define
\[
g(\tau^F,\tau^C)
=
\left( \frac{\tau^F - \bar{\tau}^F}{\sigma_F} \right)^2
+
\left( \frac{\tau^C - \bar{\tau}^C}{\sigma_C} \right)^2
+
\lambda \cdot w_{\text{region}}
\left|
\frac{(\tau^F - \bar{\tau}^F)(\tau^C - \bar{\tau}^C)}{\sigma_F \sigma_C}
\right|,
\]
where $\bar{\tau}^F$ and $\bar{\tau}^C$ are root-level treatment effects computed on the full training sample, $(\sigma_F,\sigma_C)$ are the corresponding root-level standard errors, $\lambda \ge 0$ controls the strength of the joint deviation term, and $w_{\text{region}} \in \{0,1\}$ is a region-specific \emph{activator} that determines whether the co-movement component is active.


Normalization by root-level standard errors anchors the splitting criterion to a global reference scale. This choice rewards splits that uncover substantively meaningful deviations from the overall policy effect and prevents depth-dependent rescaling that could otherwise amplify local noise. While alternative normalizations are possible, root-level normalization is well aligned with the diagnostic and auditing goals of the method.

After each candidate split, each resulting segment is assigned to a region based on the signs of its estimated treatment effects $(\hat{\tau}^F,\hat{\tau}^C)$. The analyst specifies a set of regions of interest $\mathcal{R}_{\text{interest}} \subseteq \{1,2,3,4\}$, and the $w_{\text{region}}$ equals one if the segment belongs to one of these regions and zero otherwise.

The co-movement term uses an absolute value to measure the \emph{strength} of joint deviation between firm and consumer effects, while the direction of the effects is governed entirely by region membership. This separation allows the algorithm to prioritize segments exhibiting outcome patterns of interest without conflating magnitude with sign.


By construction, splits that generate heterogeneity in only one outcome receive lower gain than splits that jointly differentiate firm and consumer responses within regions of interest. This ensures that the resulting segmentation reflects joint outcome structure rather than isolated variation in a single metric.

The splitting criterion is designed for tree construction and is not used for model selection.

\subsection{Model Selection and Cross-Validation}
\label{sec:model-selection}

While the joint splitting criterion governs how the tree is grown, it is not designed for selecting hyperparameters such as maximum tree depth or minimum leaf size. The splitting objective deliberately emphasizes within-sample heterogeneity in joint outcomes and may therefore favor overly complex trees in finite samples. To mitigate overfitting, we decouple tree construction from model selection and adopt a cross-validation strategy based on pseudo-outcomes \citep{wager2018estimation}.

The pseudo-outcome approach transforms the causal estimation problem into a supervised prediction problem whose expected loss is minimized when treatment effects are estimated accurately. Under randomized assignment with equal treatment probability, the pseudo-outcomes for individual $i$ are defined as
\[
\tilde{Y}^F_i = (2T_i - 1) Y^F_i, \qquad
\tilde{Y}^C_i = (2T_i - 1) Y^C_i,
\]
where the transformation $(2T_i - 1)$ maps the treatment indicator from $\{0,1\}$ to $\{-1,1\}$. Under randomization, the conditional expectation of each pseudo-outcome is proportional to the corresponding treatment effect.

For a given hyperparameter configuration, we perform $K$-fold cross-validation. In each fold, the tree is trained on the training folds using the joint splitting criterion, and predicted segment-level treatment effects $(\hat{\tau}^F_i, \hat{\tau}^C_i)$ are obtained for observations in the validation fold. Model performance is then evaluated using mean squared error on pseudo-outcomes:
\[
\mathcal{L}_k
=
\frac{1}{n_{\text{val},k}} \sum_{i \in \text{val}_k}
\left( \tilde{Y}^F_i - \frac{1}{2}\hat{\tau}^F_i \right)^2
+
\frac{1}{n_{\text{val},k}^C} \sum_{i \in \text{val}_k : Y_i^C \text{ observed}}
\left( \tilde{Y}^C_i - \frac{1}{2}\hat{\tau}^C_i \right)^2,
\]
where $n_{\text{val},k}$ is the number of observations in validation fold $k$, and $n_{\text{val},k}^C$ is the number of observations with observed consumer outcomes in that fold. The factor $1/2$ aligns the pseudo-outcomes with the treatment-effect parameterization under equal-probability randomization.

Hyperparameters are selected to minimize the average validation loss across folds. This selection criterion favors trees that produce stable and generalizable estimates of both firm-side and consumer-side treatment effects, while remaining agnostic to the region-activation structure used during tree construction.


\subsection{Algorithm Properties and Use Case}

The proposed algorithm is designed to (i) uncover heterogeneous joint treatment effects on firm and consumer outcomes, (ii) produce interpretable, rule-based segments, and (iii) allow analysts to prioritize specific outcome patterns through region-based weighting. Unlike approaches that estimate individual-level treatment effects and cluster ex post, the method integrates estimation and segmentation, accounts for estimation uncertainty during tree growth, and yields transparent segments suitable for auditing optimized marketing policies.

In the next section, we evaluate the behavior of the algorithm through a simulation study.

\section{Simulation Study}
\label{sec:simulation1}

We evaluate the proposed algorithm using a controlled simulation designed to assess its statistical behavior and diagnostic properties across a wide range of data environments. The purpose of this simulation is not to approximate any specific market or behavioral setting, but to provide a systematic stress test of the algorithm under conditions that vary in structural complexity, signal strength, sample size, and alignment between firm-side and consumer-side treatment effects.

The simulation is constructed to isolate the role of the co-movement emphasis parameter $\lambda$. By varying key properties of the data—such as sparsity, noise, prevalence of positive treatment effects, and overlap between firm and consumer signals—we examine how prioritizing specific joint outcome patterns affects the algorithm’s ability to detect regions of interest. This design allows us to characterize the trade-offs between overall accuracy and targeted recall that are central to auditing applications.

\subsection{Binary Data Generating Process}
\label{sec:binary-dgp}

We generate data using an abstract binary data generating process that provides explicit control over the structure of treatment effect heterogeneity. The process is fully specified and non-behavioral, allowing us to manipulate the statistical properties of the data independently of economic or psychological assumptions.

\textbf{Feature construction.}
Each observation is characterized by $k$ categorical variables, where categorical variable $j$ has $n_j$ possible categories. Each observation belongs to exactly one category per variable. All categorical variables are one-hot encoded, producing a binary feature vector $X \in \{0,1\}^p$, where
\[
p = \sum_{j=1}^k n_j.
\]
Each observation therefore has exactly $k$ active features and $p-k$ inactive features.

To ensure comparability across simulation runs, the total number of features is fixed at $p=60$. This is achieved by setting
\[
n_j = 60 / k \quad \text{for all } j,
\]
so that increasing $k$ increases sparsity and the depth of interactions required to isolate specific feature combinations, without changing overall dimensionality.

\vspace{6pt}
\noindent
\textbf{Baseline outcomes.}
Baseline firm-side and consumer-side outcomes are generated as linear functions of the binary features:
\[
\mu^F(X) = X^\top \beta^F, \qquad \mu^C(X) = X^\top \beta^C,
\]
where $\beta^F$ and $\beta^C$ are vectors of randomly drawn coefficients. These baseline components are independent of treatment assignment and do not induce heterogeneity in treatment effects.

\vspace{6pt}
\noindent
\textbf{Treatment effect structure.}
Treatment effect heterogeneity is introduced through \emph{activating feature combinations}. Let $\mathcal{C}$ denote the set of all possible category combinations implied by the $k$ categorical variables. From this set, we randomly select:
\begin{itemize}
    \item $m_F$ activating combinations for the firm outcome,
    \item $m_C$ activating combinations for the consumer outcome,
\end{itemize}
where $m_F = m_C$ defines the \emph{complexity} of the data generating process and ranges between 1 and 30. Higher values of complexity correspond to a larger number of distinct feature combinations that induce positive treatment effects, requiring more complex tree structures to recover.

The degree of overlap between the firm and consumer activating sets is governed by the \emph{covariance between firm and consumer signals}. Specifically, a fraction $\rho \in [0,1]$ of the consumer activating combinations are shared with the firm activating combinations. When $\rho = 0$, activating combinations are disjoint; when $\rho = 1$, the two sets coincide.

For any observation with feature vector $X$, treatment effects are defined as
\[
\tau^F(X) =
\begin{cases}
+\text{intensity}, & \text{if } X \in \mathcal{C}_F, \\
-\text{intensity}, & \text{otherwise},
\end{cases}
\qquad
\tau^C(X) =
\begin{cases}
+\text{intensity}, & \text{if } X \in \mathcal{C}_C, \\
-\text{intensity}, & \text{otherwise},
\end{cases}
\]
where $\mathcal{C}_F$ and $\mathcal{C}_C$ denote the firm and consumer activating sets, respectively.

Independent Gaussian noise is added to each treatment effect prior to outcome generation, allowing noise to attenuate or flip treatment-effect signs.

\vspace{6pt}
\noindent
\textbf{Assignment of observations to combinations.}
While activating combinations are fixed at the level of the feature space, the assignment of observations to combinations is random. The \emph{proportion of positive treatment effects} parameter controls the fraction of observations that are assigned to activating combinations. Specifically, a given fraction of observations is randomly assigned to activating combinations, with the remaining observations assigned to non-activating combinations. Conditional on this constraint, feature combinations are sampled uniformly at random. This design ensures that the prevalence of positive treatment effects is controlled independently of the complexity and overlap of activating combinations.

\vspace{6pt}
\noindent
\textbf{Outcome generation and region labels.}
Treatment assignment is randomized as $T \sim \mathrm{Bernoulli}(0.5)$. Outcomes are generated as
\[
Y^F = \mu^F(X) + \tau^F(X)\cdot T + \varepsilon^F, \qquad
Y^C = \mu^C(X) + \tau^C(X)\cdot T + \varepsilon^C,
\]
where $\varepsilon^F$ and $\varepsilon^C$ are independent noise terms.

Each observation is assigned a \emph{true region type} based on the signs of the noisy treatment effects $(\tau^F(X), \tau^C(X))$:
\begin{itemize}
    \item Region 1: $\tau^F > 0$ and $\tau^C > 0$ (win--win),
    \item Region 2: $\tau^F > 0$ and $\tau^C \leq 0$ (firm gain, consumer loss),
    \item Region 3: $\tau^F \leq 0$ and $\tau^C > 0$ (consumer gain, firm loss),
    \item Region 4: $\tau^F \leq 0$ and $\tau^C \leq 0$ (lose--lose).
\end{itemize}

These region labels are used exclusively for evaluation and are never observed by the algorithm.

\subsection{Experimental Factors}
\label{sec:binary-factors}

Each simulation run is characterized by a set of experimental factors that jointly determine the structure and difficulty of the data generating process. These factors are sampled independently across runs and control distinct aspects of treatment effect heterogeneity, signal quality, and statistical power. Together, they generate a broad spectrum of environments ranging from simple and well-powered to sparse, noisy, and highly imbalanced.

\begin{itemize}
    \item \emph{Complexity.}  
    Complexity controls the number of activating feature combinations that induce positive treatment effects for both firm-side and consumer-side outcomes. Specifically, it determines the size of the activating sets $\mathcal{C}_F$ and $\mathcal{C}_C$. As complexity increases, treatment effects depend on a larger number of distinct feature combinations, requiring deeper and more intricate tree structures to correctly isolate the relevant segments.

    \item \emph{Covariance between firm and consumer signals.}  
    Covariance determines the fraction of activating combinations shared between the firm and consumer activating sets. A covariance of zero implies that activating combinations for firm and consumer outcomes are disjoint, while a covariance of one implies complete overlap. Intermediate values generate partial overlap. This factor directly controls the extent to which firm-side and consumer-side treatment effects are driven by similar underlying feature combinations.

    \item \emph{Proportion of positive treatment effects.}  
    This factor governs the fraction of observations that are assigned to activating combinations. While the activating combinations themselves are fixed at the feature-space level, observations are randomly assigned to combinations subject to this constraint. When the proportion is low, most observations experience negative treatment effects for both outcomes; when it is high, most experience positive effects. As a result, trade-off regions—such as firm gain with consumer loss—are most prevalent at intermediate values.

    \item \emph{Sparsity.}  
    Sparsity is controlled by the number of categorical variables $k$. Because the total number of features is fixed, increasing $k$ reduces the number of categories per variable and increases the number of splits required to isolate a specific activating combination. Higher sparsity therefore increases the depth and structural complexity of the tree needed to recover the treatment effect heterogeneity.

    \item \emph{Noise.}  
    Noise controls the standard deviation of Gaussian perturbations added to the treatment effects prior to outcome generation. Higher noise reduces the signal-to-noise ratio, increases the likelihood of sign reversals, and makes correct region identification more difficult even when the underlying structure is simple.

    \item \emph{Training sample size.}  
    Sample size determines the number of observations available for estimation. The size of the training dataset is drawn uniformly on a log scale between 1{,}000 and 200{,}000 observations, with an independent test dataset of half the size. Larger samples improve estimation precision and support deeper segmentation, while smaller samples exacerbate variance and increase the risk of false negatives.
\end{itemize}

By varying these factors independently, the simulation explores a wide range of data environments in which the difficulty of detecting segments of interest arises from different sources, including structural complexity, limited overlap between firm and consumer signals, class imbalance, noise, and finite-sample uncertainty.


\subsection{Estimation and Experimental Protocol}
\label{sec:estimation-protocol}

This subsection describes how the proposed algorithm is estimated and evaluated within the simulation framework. The focus is on understanding how the co-movement emphasis parameter $\lambda$ influences the algorithm’s ability to detect segments of interest, rather than on comparing alternative estimation methods.

\subsubsection{Simulation design}

Each simulation run builds on the binary data generating process described in Sections~\ref{sec:binary-dgp} and~\ref{sec:binary-factors}. For every run, values for the experimental factors—complexity, noise, sparsity, proportion of positive treatment effects, covariance between firm and consumer signals, and training sample size—are independently sampled from their respective ranges. This procedure generates data environments that differ systematically in structural difficulty, signal strength, class imbalance, and alignment between firm-side and consumer-side treatment effects.

For each sampled configuration, we generate an independent training dataset and a held-out test dataset. Treatment assignment is randomized with equal probability. True region labels are determined by the signs of the underlying (noisy) treatment effects and are used exclusively for out-of-sample evaluation; they are never observed by the algorithm during estimation. To ensure robust coverage of the parameter space, we repeat this procedure across 10{,}000 independent simulations.

\subsubsection{Algorithm configurations}

We estimate multiple configurations of the proposed algorithm that differ only in the value of the co-movement emphasis parameter $\lambda$ and the specification of regions of interest. All configurations use the same model-selection procedure described in Section~\ref{sec:model-selection}, ensuring that differences in performance are attributable solely to the role of $\lambda$ rather than to differences in tree complexity or tuning.

We consider the following specifications:
\begin{itemize}
    \item A baseline configuration with $\lambda = 0$, in which the co-movement component is inactive and all regions are treated symmetrically.
    \item A set of region-targeted configurations with $\lambda \in \{1, 2, 3, 4, 6, 8, 10\}$, each specifying Region~2 as the sole region of interest.
\end{itemize}

For $\lambda > 0$, the co-movement term acts as an activator within the splitting objective: it contributes to the gain function only when a candidate segment is classified as belonging to Region~2 based on its estimated treatment effects in the training data. Increasing $\lambda$ strengthens this activation, placing greater emphasis on identifying segments where firm-side gains coincide with consumer-side losses.

\subsubsection{Evaluation metrics}

All algorithm configurations are evaluated on the same held-out test datasets to ensure direct comparability. Performance is summarized using a set of classification-based metrics derived from the true region labels. We report overall accuracy across all regions, as well as metrics focused specifically on Region~2, including recall and the false negative rate.

This evaluation protocol is intentionally aligned with the auditing objective of the method. While overall accuracy captures global classification performance, recall for Region~2 reflects the algorithm’s ability to surface segments of regulatory or ethical concern. By holding all other aspects of estimation fixed, this design isolates the role of $\lambda$ in shaping the trade-off between predictive accuracy and targeted sensitivity.

In the next subsection, we present the results of this comparison and discuss how different values of $\lambda$ affect this trade-off across the simulated environments.

\subsection{Results}
\label{sec:lambda-results}

\subsubsection{Accuracy--Recall Trade-off and Optimal $\lambda$}

We begin by examining how the co-movement emphasis parameter $\lambda$ affects the trade-off between overall accuracy and the ability to detect Region~2 segments, where firm-side gains coincide with consumer-side losses. Figure~\ref{fig:lambda-accuracy-recall} summarizes average performance across all simulation runs.

As $\lambda$ increases, overall accuracy declines monotonically. This pattern reflects the fact that activating the co-movement component in the splitting criterion prioritizes the identification of Region~2 segments, potentially at the expense of accurately assigning observations belonging to other regions. When $\lambda = 0$, the algorithm treats all regions symmetrically and achieves the highest overall accuracy, corresponding to a purely accuracy-driven objective.

In contrast, recall for Region~2 exhibits a non-monotonic relationship with $\lambda$. Relative to the baseline case $\lambda = 0$, recall increases sharply for small positive values of $\lambda$, indicating that explicitly emphasizing co-movement improves the algorithm’s ability to detect firm-benefiting but consumer-harming segments. Recall reaches its maximum at $\lambda = 4$ and then declines gradually as $\lambda$ increases further. This decline suggests that excessive emphasis on co-movement can lead to over-specialization, reducing generalization and increasing false positives in held-out data.

This divergence between accuracy and recall highlights the distinction between optimization and auditing objectives. If the goal were to maximize overall accuracy, the baseline configuration with $\lambda = 0$ would be preferred. However, the purpose of our method is auditing rather than prediction. In many regulatory and ethical contexts, false negatives are particularly costly: failing to detect segments in which firm-side improvements come at the expense of consumer welfare may lead to regulatory violations, reputational harm, or disproportionate impact on vulnerable populations. From this perspective, recall for Region~2 is a primary performance criterion.

Although $\lambda = 4$ achieves the highest recall, the incremental gain relative to $\lambda = 2$ is small. The increase in recall between these two values is minimal, while the associated reduction in overall accuracy is more pronounced. This pattern indicates diminishing returns to increasing $\lambda$ beyond moderate values.

Balancing these considerations, we view $\lambda = 2$ as a particularly attractive operating point. It captures most of the recall improvement relative to $\lambda = 0$, substantially reducing false negatives for Region~2, while avoiding the larger accuracy losses associated with higher values of $\lambda$. For practitioners seeking a default configuration aligned with auditing and compliance objectives—rather than pure predictive performance—we therefore recommend $\lambda = 2$ as a sensible default.

\subsubsection{Performance Across Data Generating Dimensions}
\label{sec:lambda-dimensions}

We next examine how the accuracy--recall trade-off varies across key dimensions of the data generating process for $\lambda \in \{0, 2, 4\}$\footnote{Additional figures and detailed discussion are provided in Appendix~\ref{app:lambda-dimensions}.}. Across dimensions, the recall curves for $\lambda = 2$ and $\lambda = 4$ are generally close, indicating no practically meaningful separation between these two configurations in most settings. This finding further supports $\lambda = 2$ as a robust default choice.

As complexity increases—that is, as the number of activating combinations grows—both overall accuracy and Region~2 recall decline, reflecting the increasing difficulty of recovering richer heterogeneous structures. At the same time, the recall gap between $\lambda = 0$ and $\lambda > 0$ becomes more pronounced at higher complexity levels, highlighting the value of the co-movement activation in settings where treatment-effect patterns are more intricate.

Covariance between firm and consumer signals yields a particularly informative pattern. While accuracy differences across $\lambda$ values remain relatively stable, the recall advantage of activating the co-movement term is minimal when firm and consumer activating combinations do not overlap and becomes largest when overlap is high. This indicates that the co-movement component is most effective when firm-side and consumer-side treatment effects are driven by similar underlying feature combinations, consistent with the design of the joint splitting objective.

Sample size also shapes the trade-off. As the number of training observations increases, accuracy improves for all configurations, and the accuracy gap between $\lambda = 0$ and $\lambda > 0$ becomes more visible. In contrast, the recall gap narrows with larger samples: the benefits of the co-movement activation are most pronounced in smaller samples and diminish as estimation precision improves.

Noise produces the expected degradation in both accuracy and recall. As noise increases, the signal-to-noise ratio decreases and sign reversals become more frequent, reducing the algorithm’s ability to correctly assign regions regardless of $\lambda$.

The proportion of observations assigned to activating combinations affects performance through the induced distribution of region types. Accuracy is highest at the extremes, where most observations fall into either win--win or lose--lose regions, and lowest at intermediate values where multiple region types coexist. In contrast, Region~2 recall increases with this proportion, and region-targeted configurations consistently outperform $\lambda = 0$ across most values. This suggests that targeted activation remains useful even as the overall region distribution shifts across environments.

Finally, sparsity—captured by the number of categorical variables $k$—reduces both accuracy and recall as the depth required to isolate activating combinations increases. When sparsity is low, differences across $\lambda$ values are small. As sparsity increases, however, the trade-off becomes more pronounced: both the accuracy loss from region targeting and the recall gains relative to $\lambda = 0$ grow larger. This pattern reinforces the value of targeted auditing in structurally complex settings.

\subsection{Conclusion of Simulation}

The simulation results demonstrate that activating the co-movement component in the splitting objective substantially improves recall for Region~2 segments, which is a central requirement for auditing applications. Across a wide range of simulated environments, configurations that emphasize region-specific co-movement consistently reduce false negatives relative to the accuracy-driven baseline with $\lambda = 0$.

On average, $\lambda = 2$ provides the most favorable trade-off between recall and accuracy. While larger values of $\lambda$ can yield further recall improvements in some settings, these gains are typically modest and are accompanied by more pronounced reductions in overall accuracy. Across most data-generating dimensions, $\lambda = 2$ performs comparably to higher values while maintaining better global performance.

The results also highlight that the value of activating the co-movement term depends on the structure of the underlying data. In complex or highly structured environments—where treatment-effect patterns are harder to recover—region targeting leads to meaningful improvements in recall. In simpler settings, by contrast, the incremental benefits may be smaller and do not always justify the associated accuracy loss. Taken together, these findings underscore the importance of explicitly prioritizing recall when the goal is auditing for potential consumer harm, rather than maximizing predictive accuracy.


\begin{figure}[t!]
    \centering
    \includegraphics[width=0.9\textwidth]{figures/figurefiles/lambda_comparison_accuracy.png}
    \caption*{(a) Overall accuracy as a function of $\lambda$.}
    \vspace{0.5em}
    \includegraphics[width=0.9\textwidth]{figures/figurefiles/lambda_comparison_recall_region_2.png}
    \caption*{(b) Recall for Region~2 as a function of $\lambda$.}
    \caption{Trade-off between overall accuracy and recall for Region~2 as the co-movement regularization parameter $\lambda$ varies. Increasing $\lambda$ improves recall for Region~2 up to a point, at the cost of reduced overall accuracy.}
    \label{fig:lambda-accuracy-recall}
\end{figure}


\begin{comment}
\section{Simulation Study 2: Free trial}
\label{sec:simulation1}
\subsection{Motivation and Overview}

The first simulation study focuses on a deliberately abstract data generating process, designed to isolate the statistical properties of the proposed algorithm and to understand how the co-movement term affects accuracy and recall under controlled heterogeneity. While this abstraction is useful for stress-testing the method, many practical concerns in consumer-centric auditing arise in settings with richer behavioral structure and clearer managerial interpretation.

The second simulation therefore serves two complementary goals. First, it introduces a more concrete and policy-relevant scenario—free trial design—where firm-side optimization decisions are common and consumer welfare considerations are central. Second, and more importantly, this simulation highlights a distinct source of consumer harm that is not driven by negative treatment effects per se, but by selection: firm-side optimization increases conversion by disproportionately attracting or retaining consumers for whom the product is not genuinely beneficial. Even when treatment effects on consumer outcomes are weak or ambiguous, such selection mechanisms can generate systematic consumer harm.

By modeling these behavioral distortions explicitly, this simulation demonstrates that the proposed auditing framework is not limited to detecting harmful treatment effects. Instead, it can also surface segments where firm-side gains arise from psychologically driven selection effects, enabling the identification of consumer harm even when the underlying issue is not a directly negative treatment effect.

\subsection{Simulation Setup}
\label{sec:simulation2-setup}

We simulate a free-trial subscription environment in which consumers are randomly assigned to either a short trial ($T=0$) or a long trial ($T=1$). Treatment assignment follows a Bernoulli distribution with equal probability. After completing the trial, each consumer decides whether to subscribe. Firm outcomes are observed for all users, while consumer outcomes are observed only for subscribers.

Each consumer $i$ is associated with a vector of observed features $X_i \in \mathbb{R}^p$, drawn independently from a standard normal distribution. These features determine latent psychological constructs through deterministic but unobserved gating mechanisms. Neither the firm nor the auditing algorithm observes these latent variables; they are introduced solely through their effect on subscription behavior.

The firm evaluates trial policies using subscription rates, while consumer welfare is evaluated using realized satisfaction conditional on subscription. The simulation generates environments in which firm-side optimization can improve conversion by selecting consumers whose subscription decisions are driven by psychological factors rather than genuine product usefulness.


\subsection{Latent Psychological Constructs}
\label{sec:simulation2-latent}

Each consumer is characterized by three latent psychological constructs: usefulness ($U_i$), novelty ($N_i$), and sunk cost ($S_i$). These constructs govern subscription behavior but are not directly observed.

Each construct is activated through a binary gate that depends on a subset of observed features. For a given construct, the gate takes the form
\[
g_i = \mathbbm{1}\left\{ X_{i,j} > \theta_j \ \text{for all } j \in \mathcal{G} \right\},
\]
where $\mathcal{G}$ denotes the indices of the features controlling the gate and $\theta_j$ are fixed thresholds. The gate equals one only when all associated features exceed their thresholds.

Latent constructs are generated using a bipolar transformation:
\[
k_i = (2g_i - 1)\cdot m + v \cdot \varepsilon_i, \qquad \varepsilon_i \sim \mathcal{N}(0,1),
\]
\[
L_i = \sigma(k_i),
\]
where $m$ controls the degree of bipolarity, $v$ controls individual-level noise, and $\sigma(\cdot)$ is the logistic function. Larger values of $m$ induce sharply separated high and low values, generating clear behavioral heterogeneity.

Usefulness represents the true long-run value of the product to the consumer and serves as the benchmark for consumer welfare. Novelty captures short-term excitement associated with trying a new product, while sunk cost reflects psychological commitment accumulated through time and effort during the trial. Novelty and sunk cost influence subscription behavior but do not directly increase realized consumer welfare.

\subsection{Perceived Value and Subscription Decision}
\label{sec:simulation2-decision}

Consumers base their subscription decision on perceived value rather than true usefulness. Perceived value depends on trial length and latent psychological constructs.

For short trials ($T=0$), perceived value is given by
\[
P_i^{\text{short}} = \log\!\left( \exp(N_i) + \exp(a_u U_i) \right),
\]
while for long trials ($T=1$),
\[
P_i^{\text{long}} = \log\!\left( \exp(U_i) + \exp(a_s S_i) \right),
\]
where $a_u$ and $a_s$ control the relative importance of usefulness in short trials and sunk cost in long trials.

Subscription probability is a linear transformation of perceived value:
\[
\pi_i = \mathrm{clip}(a_0 + a_p P_i, 0, 1),
\]
and the firm outcome is generated as
\[
Y_i^F \sim \mathrm{Bernoulli}(\pi_i).
\]

This structure implies that short trials primarily attract consumers with high novelty, while long trials retain consumers with high sunk cost, even when usefulness is low.

\subsection{Consumer Outcomes and Selection}
\label{sec:simulation2-outcomes}

Consumer satisfaction is observed only for subscribers and equals realized usefulness:
\[
Y_i^C =
\begin{cases}
U_i, & \text{if } Y_i^F = 1, \\
\text{unobserved}, & \text{otherwise}.
\end{cases}
\]

Trial length does not directly affect usefulness. Any differences in consumer outcomes across trial policies arise exclusively through selection into subscription. As a result, firm-side optimization can increase subscription rates by disproportionately selecting consumers whose perceived value is inflated by novelty or sunk cost, even when their realized satisfaction is low.

\subsection{Results}
\label{sec:simulation2-results}

The results show that the auditing algorithm identifies segments where long trials increase subscription rates while lowering average consumer welfare among subscribers. These segments correspond to populations for which sunk cost or novelty dominates usefulness in perceived value formation.

Despite not observing the latent behavioral constructs, the algorithm recovers partitions that align with the underlying selection mechanisms. In environments where behavioral distortions are weak, the method refrains from producing splits, indicating robustness against spurious findings.

Overall, the detected divergence between firm and consumer outcomes is driven by endogenous selection rather than by negative treatment effects on consumer outcomes.

\subsection{Conclusion of Simulation~2}

This simulation demonstrates that consumer harm can arise from firm-side optimization through behavioral selection mechanisms, even when treatment effects on consumer outcomes are not directly negative. By jointly auditing firm and consumer outcomes, the proposed framework identifies segments where firm gains are achieved by selecting consumers for whom the product is not genuinely beneficial.

Together with Simulation~1, these results show that the framework applies to both treatment-effect-driven and selection-driven forms of consumer harm, broadening its relevance for auditing data-driven optimization policies.

\end{comment}

\section{Conclusions}
\label{sec:conclusion}

Marketing interventions are increasingly selected through continuous experimentation and optimization against firm-side metrics such as conversion, engagement, and revenue. While this infrastructure has improved managerial discipline and learning speed, it also creates a practical accountability problem: when what is optimized is narrowly defined and continuously updated, consumer-side consequences can remain unmeasured, delayed, or visible only in aggregate. In such environments, harms are most likely to be localized—concentrated in specific segments—while overall firm gains appear unambiguously positive. This gap between optimization and accountability is likely to widen further as generative AI reduces the cost of producing and iterating persuasive, personalized content at scale.

This paper proposes an auditing approach that leverages the same randomized A/B testing data firms already generate. Rather than introducing a new targeting policy, the goal is diagnostic: to identify interpretable subpopulations for which an optimized intervention creates systematically different joint responses in a firm-side outcome and a consumer-side outcome. We extend recursive partitioning to a two-outcome setting and introduce a joint splitting criterion that evaluates candidate splits based on how they reveal heterogeneity in both outcomes, not only in one. In addition, a region-based weighting component allows the audit to emphasize outcome patterns of substantive interest—such as segments where firm-side gains coincide with consumer-side losses—without conflating the sign pattern with the magnitude of deviation. The resulting tree provides transparent, rule-based segments that can be communicated to managers, compliance teams, and researchers as segment-level diagnostics of alignment versus divergence.

The simulation study highlights a central trade-off that is inherent to auditing applications. When the objective is to maximize overall classification accuracy across all joint-outcome regions, a symmetric criterion (i.e., without region emphasis) performs best. However, when the purpose is to detect segments of concern—particularly firm-win/consumer-lose segments—false negatives can be especially costly. Activating the region-weighted co-movement component increases sensitivity to these segments and improves recall, with moderate values of the emphasis parameter achieving substantial reductions in false negatives while limiting the loss in global accuracy. These results illustrate that auditing differs from prediction: the appropriate operating point depends on the relative cost of missed detections versus false alarms. Importantly, our estimation procedure is designed to mitigate overfitting and enhance stability. We decouple tree construction from model selection, select complexity using cross-validation on pseudo-outcomes, and use pruning to produce stable, parsimonious trees that generalize to held-out data and remain interpretable for auditing and governance purposes.

The proposed framework has several implications for marketing practice and research. For firms, the method provides a practical tool for responsible personalization and internal governance. Because it uses standard experimental data and produces interpretable segmentation rules, it can be embedded into existing experimentation pipelines as a recurring diagnostic: after an intervention is chosen or updated, the audit can reveal whether firm-side improvements are accompanied by localized consumer-side losses, and which segments are most affected. This supports actionable responses such as imposing segment-level guardrails, triggering additional measurement (e.g., targeted satisfaction or regret surveys), altering the intervention for flagged segments, or redesigning objective functions to incorporate consumer outcomes. For researchers, the method offers a nonparametric way to empirically characterize when marketing optimization aligns with consumer outcomes and when it diverges, beyond pricing settings and beyond average effects. It enables systematic study of the conditions under which outcome misalignment emerges (e.g., by user history, engagement intensity, or contextual variables), and it can be used to evaluate how new optimization technologies—such as LLM-driven content generation—shift the distribution of benefits and costs across consumers.

More broadly, the approach provides a bridge between optimization-focused marketing analytics and the growing emphasis on accountability in digital markets. Because marketing policies are adaptive and continuously re-optimized, one-time evaluations are insufficient. Segment-level auditing based on randomized experimentation data offers a scalable and repeatable way to monitor downstream consequences over time, supporting both internal oversight and external scrutiny in environments where interventions evolve rapidly. At the same time, the framework is intentionally modular: the “consumer-side outcome” can be defined in different ways depending on context, ranging from usage and retention to satisfaction, regret, returns, or other measures of welfare that are meaningful for the domain.

Several extensions are promising. First, many applications involve more than two outcomes (e.g., short-run conversion, long-run retention, satisfaction, and trust), motivating multi-outcome auditing and segmentation based on richer notions of alignment, divergence, and Pareto trade-offs. Second, many consumer consequences unfold over time; repeated audits and longitudinal outcome measurement could help distinguish transient effects from persistent welfare changes and clarify the dynamics of optimization-induced harm. Third, when consumer outcomes are partially observed due to post-treatment selection, additional methodological work on sensitivity analysis and robust estimation would further strengthen the audit’s interpretability and credibility. Finally, an important managerial direction is to connect diagnosis to redesign: audit results can inform constrained optimization, segmented policy adjustments, or multi-objective experimentation, followed by re-auditing in an iterative governance loop.

Overall, the main message of the paper is that experimentation infrastructure can serve not only optimization, but also accountability. By turning otherwise implicit firm--consumer trade-offs into measurable, interpretable, segment-level diagnostics, the proposed method provides a practical way to evaluate whether optimized marketing policies help or harm different consumers. This, in turn, offers a path toward personalization systems that are not only effective for firms, but also governable and aligned with marketing’s foundational promise of shared value creation.



\newpage
\renewcommand{\baselinestretch}{1.0} 
\putbib


\end{bibunit}
\begin{bibunit}

\newpage

\setcounter{table}{0}
\setcounter{page}{1}
\setcounter{figure}{0}
\setcounter{equation}{0}
\setcounter{footnote}{0}
\renewcommand{\thetable}{A\arabic{table}}
\renewcommand{\thefigure}{A\arabic{figure}}
\renewcommand{\theequation}{A.\arabic{equation}}
\renewcommand{\thepage}{\roman{page}}
\title{Are We Helping or Harming? A Method to Audit Optimized Marketing Policies}


\author{
       Ebrahim Barzegary  \\
        ESSEC Business school\\
 }

\maketitle
\begin{appendices}

\section{Additional Results: Variation Across Data Generating Dimensions}
\label{app:lambda-dimensions}

This appendix reports how the performance of the algorithm varies across key dimensions of the data generating process. For each dimension, we bin simulation runs by the factor value and plot average performance for $\lambda\in\{0,2,4\}$. The primary goal is to understand when region-weighted splitting (Region~2 as the region of interest) provides the largest gains in Region~2 recall and how those gains trade off against overall accuracy.

Across dimensions, $\lambda=2$ and $\lambda=4$ exhibit similar performance, and differences between them are typically small relative to the gap between $\lambda=0$ and $\lambda>0$. This reinforces the practical recommendation of $\lambda=2$ as a default configuration.

\subsection{Complexity}
As complexity increases (more activating combinations), both overall accuracy and Region~2 recall decrease for all methods, reflecting a more difficult heterogeneity structure. At the same time, the recall gap between $\lambda=0$ and $\lambda\in\{2,4\}$ grows under higher complexity, suggesting that region-weighted splitting is particularly valuable in settings where treatment-effect patterns are more complex and require more structure to recover.

\begin{figure}[t!]
    \centering
    \includegraphics[width=0.92\textwidth]{Figures/FiguresFile/factor_comparison_complexity_accuracy.png}
    \caption*{(a) Accuracy across complexity.}
    \vspace{0.5em}
    \includegraphics[width=0.92\textwidth]{Figures/FiguresFile/factor_comparison_complexity_recall_region_2.png}
    \caption*{(b) Recall for Region~2 across complexity.}
    \caption{Performance across complexity (number of activating combinations). Higher complexity reduces both accuracy and Region~2 recall, while increasing the recall gap between $\lambda=0$ and $\lambda>0$.}
    \label{fig:lambda-complexity}
\end{figure}

\subsection{Covariance between firm and consumer signals}
Accuracy differences across $\lambda$ are broadly stable across covariance levels. However, the recall advantage of region-weighting depends strongly on covariance: when overlap between firm and consumer activating combinations is near zero, the recall gap between $\lambda=0$ and $\lambda\in\{2,4\}$ is small; when overlap is high, region-weighting yields substantially higher Region~2 recall. This pattern highlights that the co-movement term is most effective when firm- and consumer-side effects share underlying drivers, allowing the algorithm to exploit joint structure to identify Region~2 segments.

\begin{figure}[t!]
    \centering
    \includegraphics[width=0.92\textwidth]{Figures/FiguresFile/factor_comparison_covariance_accuracy.png}
    \caption*{(a) Accuracy across covariance.}
    \vspace{0.5em}
    \includegraphics[width=0.92\textwidth]{Figures/FiguresFile/factor_comparison_covariance_recall_region_2.png}
    \caption*{(b) Recall for Region~2 across covariance.}
    \caption{Performance across covariance (proportion of shared activating combinations). The recall gain from region-weighting is smallest when overlap is low and largest when overlap is high.}
    \label{fig:lambda-covariance}
\end{figure}

\subsection{Training sample size}
As training sample size increases, overall accuracy increases for all methods, as expected from improved estimation precision. The accuracy gap between $\lambda=0$ and $\lambda>0$ becomes more visible with larger samples. In contrast, the recall gap shrinks as sample size grows: region-weighting provides the largest improvements in Region~2 recall when sample size is relatively small, and the benefit fades as all methods become sufficiently well-powered to detect Region~2 segments.

\begin{figure}[t!]
    \centering
    \includegraphics[width=0.92\textwidth]{Figures/FiguresFile/factor_comparison_data_size_accuracy.png}
    \caption*{(a) Accuracy across training sample size.}
    \vspace{0.5em}
    \includegraphics[width=0.92\textwidth]{Figures/FiguresFile/factor_comparison_data_size_recall_region_2.png}
    \caption*{(b) Recall for Region~2 across training sample size.}
    \caption{Performance across training sample size. As sample size increases, accuracy improves for all methods and the recall gap between $\lambda=0$ and $\lambda>0$ diminishes.}
    \label{fig:lambda-datasize}
\end{figure}

\subsection{Noise}
Noise degrades both overall accuracy and Region~2 recall across all $\lambda$ values. As effect noise increases, signal-to-noise ratio decreases and sign reversals become more frequent, making region identification more difficult. Region-weighting remains beneficial for Region~2 recall at moderate noise levels, but extreme noise eventually dominates and performance declines for all methods.

\begin{figure}[t!]
    \centering
    \includegraphics[width=0.92\textwidth]{Figures/FiguresFile/factor_comparison_noise_accuracy.png}
    \caption*{(a) Accuracy across noise.}
    \vspace{0.5em}
    \includegraphics[width=0.92\textwidth]{Figures/FiguresFile/factor_comparison_noise_recall_region_2.png}
    \caption*{(b) Recall for Region~2 across noise.}
    \caption{Performance across effect noise. Both accuracy and Region~2 recall decrease as noise increases.}
    \label{fig:lambda-noise}
\end{figure}

\subsection{Proportion in Activating Combinations (Positive-Rate)}
This factor controls the proportion of observations assigned to activating combinations (i.e., those with positive treatment effects by construction). It directly shapes the distribution of region types. Overall accuracy is highest at the extremes because region-type variation decreases: when the positive-rate is very low, most observations fall into a lose--lose regime, whereas when it is very high, most fall into a win--win regime. Accuracy is lowest at intermediate values where region labels are more mixed and classification is more difficult. Region~2 recall increases with the positive-rate, and region-weighted specifications ($\lambda=2$ and $\lambda=4$) consistently achieve higher recall than the baseline $\lambda=0$, while remaining close to each other.

\begin{figure}[t!]
    \centering
    \includegraphics[width=0.92\textwidth]{Figures/FiguresFile/factor_comparison_rareness_accuracy.png}
    \caption*{(a) Accuracy across the proportion in activating combinations (positive-rate).}
    \vspace{0.5em}
    \includegraphics[width=0.92\textwidth]{Figures/FiguresFile/factor_comparison_rareness_recall_region_2.png}
    \caption*{(b) Recall for Region~2 across the proportion in activating combinations.}
    \caption{Performance across the proportion in activating combinations. Accuracy is U-shaped due to changes in region-type variation, while Region~2 recall increases with the positive-rate and improves under region-weighted objectives.}
    \label{fig:lambda-positive-rate}
\end{figure}




\subsection{Sparsity}
Sparsity (the number of categorical variables $k$) governs the depth required to isolate activating combinations. As $k$ increases, both accuracy and Region~2 recall decrease for all methods. Differences across $\lambda$ are small when sparsity is low, but become more pronounced as sparsity increases: the accuracy cost of region-weighting increases, while the recall advantage relative to $\lambda=0$ becomes more meaningful. This pattern suggests that targeted auditing is most useful when the segmentation structure requires deeper trees and is harder to recover.

\begin{figure}[t!]
    \centering
    \includegraphics[width=0.92\textwidth]{Figures/FiguresFile/factor_comparison_sparsity_accuracy.png}
    \caption*{(a) Accuracy across sparsity.}
    \vspace{0.5em}
    \includegraphics[width=0.92\textwidth]{Figures/FiguresFile/factor_comparison_sparsity_recall_region_2.png}
    \caption*{(b) Recall for Region~2 across sparsity.}
    \caption{Performance across sparsity (number of categorical variables $k$). Higher sparsity reduces both accuracy and Region~2 recall; the trade-off across $\lambda$ becomes more salient as $k$ increases.}
    \label{fig:lambda-sparsity}
\end{figure}
\begin{comment}
    

\newpage
\putbib
\thispagestyle{empty}

\end{appendices}
\end{bibunit}
\end{comment}

\end{document}
